% =============================================================================
% Standalone compilable draft — Lettuce Yield Prediction (Segmentation Sections)
% Compile:  pdflatex segmentation_paper && bibtex segmentation_paper && pdflatex segmentation_paper && pdflatex segmentation_paper
% Or use latexmk: latexmk -pdf segmentation_paper.tex
% =============================================================================

\documentclass[11pt,a4paper,twocolumn]{article}

\usepackage[margin=1in]{geometry}
\usepackage{amsmath,amssymb}
\usepackage{booktabs}
\usepackage{graphicx}
\usepackage{hyperref}
\usepackage{microtype}
\usepackage{enumitem}
\usepackage{authblk}

% Two-column tuning (optional)
\setlength{\columnsep}{0.22in}

\hypersetup{
    colorlinks=true,
    linkcolor=blue,
    citecolor=blue,
    urlcolor=blue
}

\title{Comparative Segmentation and Feature Extraction for\\
Hydroponic Lettuce Fresh-Weight Prediction in Grow-Tent Timelapse Imagery}

\author[1]{Your Name}
\author[1]{Co-Author}
\affil[1]{Department, University, Country}
\date{\today}

\begin{document}

\maketitle

\begin{abstract}
We compare three literature-backed segmentation approaches---ExGR with Otsu
thresholding, the Segment Anything Model (SAM), and FastSAM---on matched
grow-tent timelapse frames from two hydroponic lettuce trials. All methods share
identical post-processing for fair evaluation. On 53 trial-day frames,
FastSAM achieves the highest canopy coverage (12.9\%), near-SAM plant-detection
counts, and strong FastSAM--SAM mask agreement (Dice $= 0.603$) at approximately
three times lower latency than SAM. ExGR is fastest but captures only 6.1\%
canopy coverage. Based on quantitative and qualitative evidence, FastSAM is
selected for downstream fresh-weight feature extraction and yield modelling.
\end{abstract}

\noindent\textbf{Keywords:} lettuce yield prediction; hydroponic agriculture;
instance segmentation; FastSAM; SAM; ExGR; timelapse phenotyping.

% =============================================================================
% Sections 2--4: Related Work, Segmentation Methods, Results & Selection
% Include in your main document via: \input{paper/sections_segmentation}
% =============================================================================

\section{Related Work}
\label{sec:related}

\noindent
Computer-vision pipelines for hydroponic lettuce monitoring commonly rely on an image
segmentation step to extract canopy shape and size features for downstream biomass or fresh
weight modelling~\cite{gao2022}. Classical vegetation-index methods (e.g., ExG/ExGR) provide a
fast, training-free baseline for plant/background separation in controlled environments; Meyer
and Neto~\cite{meyer2008} showed ExGR (ExG$-$ExR) can be effective for greenhouse imagery, and
Otsu thresholding~\cite{otsu1979} is widely used to automate binarisation under variable
illumination. Recently, foundation models enable zero-shot segmentation: SAM~\cite{kirillov2023}
achieves strong general-purpose masks, while FastSAM~\cite{zhao2023fastsam} targets similar
quality with lower latency. A greenhouse lettuce study reported SAM and FastSAM are comparable
in accuracy but differ substantially in runtime~\cite{cornell2024}. Motivated by this, we
compare ExGR+Otsu, SAM, and FastSAM on the same grow-tent timelapse frames under identical
post-processing, and select one segmentation method for the yield-prediction pipeline.


\section{Segmentation Methods}
\label{sec:methods}

We evaluate three methods that require no training on our dataset: (i) ExGR+Otsu thresholding,
(ii) SAM zero-shot, and (iii) FastSAM zero-shot. Each method outputs candidate instance masks
per image. For fair comparison, all masks are filtered with the same size limits and the same
plant/pipe rejection heuristics (Section~\ref{sec:methods_postprocess}).

\subsection{Method A: ExGR with Otsu Thresholding and Connected Components}
\label{sec:methods_exgr}

Given an RGB image, channels are normalised to $[0,1]$ and the Excess Green and Excess Red
indices are computed as
\begin{align}
    \mathrm{ExG}  &= 2G' - R' - B', \\
    \mathrm{ExR}  &= 1.4R' - B', \\
    \mathrm{ExGR} &= \mathrm{ExG} - \mathrm{ExR}.
    \label{eq:exgr}
\end{align}
The ExGR map is binarised using Otsu's threshold~\cite{otsu1979}, then refined using simple
morphological operations. Connected components are treated as instances, and small blobs are
discarded. This baseline is fast and training-free.

\subsection{Method B: SAM (Segment Anything Model, Zero-Shot)}
\label{sec:methods_sam}

SAM is a promptable foundation model pre-trained on large-scale segmentation data.
We use the Ultralytics implementation with the \texttt{sam\_b.pt} checkpoint in automatic
mask-generation mode (no fine-tuning on our lettuce images), consistent with zero-shot
protocols in recent agricultural phenotyping work~\cite{cornell2024}.
For each full-resolution timelapse frame ($4056 \times 3040$ pixels), SAM proposes a set of
instance masks; masks are resized to the image grid, thresholded at 0.5, and filtered by area
(8{,}000--1{,}200{,}000 pixels).
SAM is typically accurate but relatively slow for large batch processing~\cite{cornell2024}.

\subsection{Method C: FastSAM (Fast Segment Anything Model)}
\label{sec:methods_fastsam}

FastSAM replaces SAM's transformer backbone with a faster CNN segmenter while retaining
prompt-free, multi-instance behaviour~\cite{zhao2023fastsam}.
We use the \textbf{FastSAM-x} checkpoint via Ultralytics with confidence 0.25, IoU 0.7, and
\texttt{retina\_masks=True}.
Inference resolution defaults to 1024\,px with automatic fallback to smaller sizes on GPU
memory limits.
FastSAM is designed for high-throughput settings where many regions must be segmented per
frame.

\subsection{Common Post-Processing and Feature Pathway}
\label{sec:methods_postprocess}

Regardless of segmentation method, each candidate mask undergoes the same validation:
size filtering, HSV-based ``plant-colour'' checks, and spatial heuristics to reject common pipe
regions. Accepted masks are saved as RGBA crops for feature extraction.

\medskip

For the \textbf{production yield pipeline}, only FastSAM outputs were carried forward at full
scale through the following stages:
\begin{enumerate}
    \item \textbf{QA filtering} --- blur, colour, and size scoring of candidate crops;
    \item \textbf{Cup assignment} --- $k$-means clustering on crop centroids ($\sim$14 cups per tray);
    \item \textbf{Daily cup features} --- selection of the best crop per cup per calendar day;
    \item \textbf{Checkpoint windows} --- aggregation over a four-day lookback before each harvest day;
    \item \textbf{ML export} --- modelling tables under \texttt{ml/}.
\end{enumerate}
The primary modelling file is \texttt{checkpoint\_trial1\_image.csv}, containing eight labelled
harvest checkpoints with image-derived window features and growth slopes; sensor features are
merged in companion tables for combined models.


\section{Segmentation Comparison and Selection for Yield Prediction}
\label{sec:results_selection}

\subsection{Experimental Protocol}
\label{sec:protocol}

We compared all three methods on \textbf{53 matched frames}: one representative image per
trial-day (Trial~1 and Trial~2), selected at the hour closest to midday for consistent
lighting.
Each method segmented the same source files; runtime, mask counts, canopy coverage, and
pairwise overlap were recorded.
Dice and IoU are computed between methods on the union of accepted plant masks, measuring
\emph{inter-method agreement} (not ground-truth accuracy).

\subsection{Quantitative Results}
\label{sec:quant_results}

Table~\ref{tab:seg_comparison} summarises aggregate performance across 53 matched frames.
FastSAM achieves the highest canopy coverage (12.9\%) with near-SAM plant-crop counts, while
being about $3\times$ faster than SAM. ExGR is fastest but captures substantially less canopy.

\begin{table}[t]
    \centering
    \caption{Segmentation comparison on 53 matched trial-day frames (aggregate means).}
    \label{tab:seg_comparison}
    \small
    \resizebox{\columnwidth}{!}{%
    \begin{tabular}{@{}lrrrrrr@{}}
        \toprule
        \textbf{Method} & \textbf{Time (s)} & \textbf{Raw} & \textbf{Crops} &
        \textbf{Cov.} & \textbf{Colour} & \textbf{Score} \\
        \midrule
        \textbf{FastSAM}  & \textbf{1.53} & 32.0 & \textbf{21.0} & \textbf{0.129} & 0.641 & \textbf{0.747} \\
        ExGR + Otsu       & \textbf{0.60} & 23.0 & 17.1 & 0.061 & \textbf{0.711} & 0.691 \\
        SAM               & 4.55 & 31.6 & \textbf{21.7} & 0.112 & 0.671 & 0.663 \\
        \bottomrule
        \multicolumn{7}{@{}l@{}}{\footnotesize Score combines crop-count proximity, coverage, colour, and speed (details in code).} \\
    \end{tabular}%
    }
\end{table}

FastSAM and SAM agree more strongly with each other than either does with ExGR (mean Dice
0.603 and IoU 0.452 for FastSAM vs.\ SAM on plant unions).

At full timelapse scale, the FastSAM pipeline processed 821 frames, producing 16{,}980
candidate crops; 3{,}637 (21.4\%) passed QA filtering; and 243 cup-day feature rows were
obtained after cup assignment and daily aggregation.

\subsection{Discussion and Method Selection}
\label{sec:discussion_selection}

ExGR is the fastest method and yields the highest per-mask plant-colour score, but its canopy
coverage is roughly half that of FastSAM, and its overlap with both foundation models is poor
($\mathrm{IoU} < 0.11$).
This indicates that classical index thresholding alone is insufficient for our multi-cup
grow-tent geometry.

SAM achieves the highest raw plant-crop count and strong overlap with FastSAM, supporting the
finding of Cornell~\cite{cornell2024} that the two models perform similarly on dense lettuce
imagery.
However, SAM requires approximately three times longer inference (4.55\,s vs.\ 1.53\,s per
frame), which is prohibitive for batch processing of more than 800 timelapse images.

\textbf{FastSAM} offers the best practical trade-off for our dataset: highest vegetation
coverage (12.9\%), near-SAM plant counts, strong FastSAM--SAM agreement, and acceptable
throughput. We therefore \textbf{selected FastSAM for downstream yield prediction}.

\subsection{Link to the Yield-Modelling Pipeline}
\label{sec:yield_pipeline}

Segmentation outputs feed the yield study by (i) extracting daily per-cup image features, (ii)
merging with daily sensors, and (iii) forming 4-day checkpoint windows aligned to harvest-day
fresh-weight labels. Only FastSAM-derived feature tables are used for yield-model training.


\bibliographystyle{plain}
\bibliography{references_segmentation}

\end{document}
